%! Author = lili
%! Date = 8/20/26


\newglossaryentry{exon-shuffling}{
    name={exon shuffling},
    description={Die Hypothese, das Gene durch
    Rekombination von verschiedenen Exons evolviert
    sind, und die frühen Exons für einzelne funktionelle
    Protein-Domänen kodieren.},
    type=defi
}

\newglossaryentry{introns-first}{
    name={“introns first” hypothesis},
    description={Gene für Proteine
    oder für (katalytisch aktive) RNAs sind
    wahrscheinlich in einer Form mit Introns
    entstanden.},
    type=defi
}

\newglossaryentry{introns-late}{
    name={“introns late” model},
    description={Die frühen Gene enthielten
    keine Introns, sondern die Introns sind erst später
    in der Evolution in einige Gene aufgenommen
    worden.},
    type=defi
}


\newglossaryentry{zform}{
    name={z-Form},
    description={Linksgängige Doppelhelix},
    type=defi
}

\newglossaryentry{umwelteinfluss}{
    name={Umwelteinfluss auf den Phänotyp},
    description={Modifikation der Ausprägung eines Genotyps durch äußere Umweltfaktoren, z. B. Temperatur oder Ernährung, sodass derselbe Genotyp unterschiedliche Phänotypen hervorbringen kann.},
    type=defi
}

\newglossaryentry{myb}{
    name={MYB},
    description={Familie von Transkriptionsfaktoren mit charakteristischer MYB-DNA-Bindedomäne, reguliert unter anderem Zelldifferenzierung und Sekundärstoffwechsel in Pflanzen.},
    type=defi
}


\newglossaryentry{dihydrouridine}{
    name={Dihydrouridin},
    description={Modifiziertes Uridin mit gesättigter C5--C6-Doppelbindung, vorwiegend in der D-Schleife der tRNA.
    Erhöht die Flexibilität der tRNA durch Destabilisierung der Stapelwechselwirkungen.},
    type=defi
}


\newglossaryentry{bhlh}{
    name={bHLH Domäne},
    description={DNA-Bindedomäne von Transkriptionsfaktoren, bestehend aus einer basischen Region zur DNA-Bindung und zwei Alpha-Helices, die die Dimerisierung mit einem weiteren bHLH-Protein ermöglichen.},
    type=defi
}

\newglossaryentry{bzip}{
    name={bZIP Domäne},
    description={ DNA-Bindedomäne von Transkriptionsfaktoren, bestehend aus einer basischen Region zur DNA-Bindung und einem Leucin-Zipper-Motiv, das die Dimerisierung zweier Untereinheiten vermittelt.},
    type=defi
}



\newglossaryentry{bform}{
    name={b-Form},
    description={Rechtsgängige doppelsträngige DNA-Doppelhelix, die im Vergleich zur A-Form entspannter ist und deren
    Nukleinbasen orthogonal zur Helixachse ausgerichtet sind. Sie ist die unter physiologischen Bedingungen typisch
    vorliegende Form doppelsträngiger DNA, z. B. bei Chromosomen in Zellen},
    type=defi
}

\newglossaryentry{aform}{
    name={a-Form},
    description={Rechtsgängige doppelsträngige DNA-Doppelhelix, die im Vergleich zur B-Form kompakter ist und deren
    Nukleinbasen nicht orthogonal zur Helixachse ausgerichtet sind. Sie ist die bei niedriger Feuchtigkeit vorliegende
    Form doppelsträngiger DNA, z. B. bei getrockneter DNA oder DNA-Kristallen},
    type=defi
}
